\par\medskip
\hypertarget{native-074-article056}{}\pdfbookmark[2]{56.}{article-056}
\begin{center}56.\end{center}
Praecepta praecedentia exemplo illustravisse haud inutile erit. Sit $\AscendingNode=172^0\,28^{\prime}\,13^{\prime\prime}{,}7$, $i=34^0\,38^{\prime}\,1^{\prime\prime}{,}1$; porro sit planum novum aequatori parallelum, adeoque $n=180^0$; angulum $\epsilon$, qui erit obliquitas eclipticae, statuimus $=23^0\,27^{\prime}\,55^{\prime\prime}{,}8$. Habemus itaque
\begin{center}
\begin{tabular}{@{}ll@{}}
$\AscendingNode-n=$ & $-7^0\,31^{\prime}\,46^{\prime\prime}{,}3$\\[.3em]
$i+\epsilon=$ & $58\quad5\quad56{,}9$\\[.3em]
$i-\epsilon=$ & $11\quad10\quad5{,}3$\\[.3em]
\makebox[18em][l]{$\log\sin\tfrac{1}{2}(\AscendingNode-n)$\hfill .\,.} & $8{,}817\,3026_n$\\[.3em]
\makebox[18em][l]{$\log\sin\tfrac{1}{2}(i+\epsilon)$\hfill .\,.\,.} & $9{,}686\,2484$\\[.3em]
\makebox[18em][l]{$\log\cos\tfrac{1}{2}(i+\epsilon)$\hfill .\,.\,.} & $9{,}941\,6108$
\end{tabular}
\end{center}
\begin{center}
\begin{tabular}{@{}ll@{}}
$\tfrac{1}{2}(\AscendingNode-n)=$ & $-3^0\,45^{\prime}\,53^{\prime\prime}{,}15$\\[.3em]
$\tfrac{1}{2}(i+\epsilon)=$ & $29\quad2\quad58{,}45$\\[.3em]
$\tfrac{1}{2}(i-\epsilon)=$ & $5\quad35\quad2{,}65$\\[.3em]
\makebox[18em][l]{$\log\cos\tfrac{1}{2}(\AscendingNode-n)$\hfill .\,.\,.} & $9{,}999\,0618$\\[.3em]
\makebox[18em][l]{$\log\sin\tfrac{1}{2}(i-\epsilon)$\hfill .\,.\,.\,.} & $8{,}988\,1405$\\[.3em]
\makebox[18em][l]{$\log\cos\tfrac{1}{2}(i-\epsilon)$\hfill .\,.\,.\,.} & $9{,}997\,9342$
\end{tabular}
\end{center}
\noindent Hinc fit
\begin{center}
\begin{tabular}{@{}ll@{}}
\makebox[18em][l]{$\log\sin\tfrac{1}{2}i'\sin\tfrac{1}{2}(\AscendingNode'+\Delta)$\hfill .\,.\,.} & $8{,}503\,5510_n$\\[.3em]
\makebox[18em][l]{$\log\sin\tfrac{1}{2}i'\cos\tfrac{1}{2}(\AscendingNode'+\Delta)$\hfill .\,.\,.} & $8{,}987\,2023$\\[.3em]
\hline
unde $\tfrac{1}{2}(\AscendingNode'+\Delta)=$ & $341^0\,49^{\prime}\,19^{\prime\prime}{,}01$\\[.3em]
\makebox[18em][l]{$\log\sin\tfrac{1}{2}i'$\hfill .\,.\,.\,.\,.\,.\,.\,.\,.\,.\,.\,.\,.} & $9{,}009\,4368$
\end{tabular}
\end{center}
\begin{center}
\begin{tabular}{@{}ll@{}}
\makebox[18em][l]{$\log\cos\tfrac{1}{2}i'\sin\tfrac{1}{2}(\AscendingNode'-\Delta)$\hfill .\,.\,.} & $8{,}758\,9134_n$\\[.3em]
\makebox[18em][l]{$\log\cos\tfrac{1}{2}i'\cos\tfrac{1}{2}(\AscendingNode'-\Delta)$\hfill .\,.\,.} & $9{,}996\,9960$\\[.3em]
\hline
unde $\tfrac{1}{2}(\AscendingNode'-\Delta)=$ & $356^0\,41^{\prime}\,31^{\prime\prime}{,}43$\\[.3em]
\makebox[18em][l]{$\log\cos\tfrac{1}{2}i'$\hfill .\,.\,.\,.\,.\,.\,.\,.\,.\,.\,.\,.\,.} & $9{,}997\,7202$
\end{tabular}
\end{center}
\noindent Obtinemus itaque $\frac{1}{2}i'=5^0\,51^{\prime}\,56^{\prime\prime}{,}445$, $i'=11^0\,43^{\prime}\,52^{\prime\prime}{,}89$,\hfil\break $\AscendingNode'=338^0\,30^{\prime}\,50^{\prime\prime}{,}44$, $\Delta=-14^0\,52^{\prime}\,12^{\prime\prime}{,}42$. Ceterum punctum $n$ in sphaera coelesti manifesto respondet- aequinoctio autumnali; quocirca distantia nodi ascendentis orbitae in aequatore ab aequinoctio vernali (eius \textit{rectascensio}) erit $158^0\,30^{\prime}\,50^{\prime\prime}{,}44$.

\SourceLeaf{075}Ad illustrationem art. 53 hoc exemplum adhuc ulterius continuabimus, formulasque pro coordinatis respectu trium planorum per Solem transeuntium evolvemus, quorum unum aequatori parallelum sit, duorumque reliquorum poli positivi in ascensione recta $0^0$ et $90^0$ sint siti: distantiae ab his planis sint resp. $z$, $x$, $y$. Iam si insuper distantia loci heliocentrici in sphaera coelesti a punctis $\AscendingNode$, $\AscendingNode'$ resp. denotetur per $u$, $u'$, fiet $u'=u-\Delta=u+14^0\,52^{\prime}\,12^{\prime\prime}{,}42$, et quae in art. 53 per $i$, $N-\AscendingNode$, $u$ exprimebantur, hic erunt $i'$, $180^0-\AscendingNode'$, $u'$. Sic per formulas illic datas prodit
\begin{center}
\begin{tabular}{@{}ll@{}}
\makebox[18em][l]{$\log a\sin A$\hfill .\,.\,.\,.} & $9{,}968\,7197_n$\\[.3em]
\makebox[18em][l]{$\log a\cos A$\hfill .\,.\,.\,.} & $9{,}554\,6380_n$\\[.3em]
\hline
unde $A=$ & $248^0\,55^{\prime}\,22^{\prime\prime}{,}97$\\[.3em]
\makebox[18em][l]{$\log a$\hfill .\,.\,.\,.\,.\,.} & $9{,}998\,7923$
\end{tabular}
\end{center}
\begin{center}
\begin{tabular}{@{}ll@{}}
\makebox[18em][l]{$\log b\sin B$\hfill .\,.\,.\,.} & $9{,}563\,8058$\\[.3em]
\makebox[18em][l]{$\log b\cos B$\hfill .\,.\,.\,.} & $9{,}959\,5519_n$\\[.3em]
\hline
unde $B=$ & $158^0\,5^{\prime}\,54^{\prime\prime}{,}97$\\[.3em]
\makebox[18em][l]{$\log b$\hfill .\,.\,.\,.\,.\,.} & $9{,}992\,0848$
\end{tabular}
\end{center}
\noindent Habemus itaque
\[
\begin{aligned}
x&=ar\sin(u'+248^0\,55^{\prime}\,22^{\prime\prime}{,}97)\\
&=ar\sin(u+263^0\,47^{\prime}\,35^{\prime\prime}{,}39)\\[.4em]
y&=br\sin(u'+158\quad5\quad54{,}97)\\
&=br\sin(u+172\quad58\quad7{,}39)\\[.4em]
z&=cr\sin u'\\
&=cr\sin(u+14\quad52\quad12{,}42),
\end{aligned}
\]
ubi $\log c=\log\sin i'=9{,}308\,1870$.

Alia solutio problematis hic tractati invenitur in \textsc{von Zach} \textit{Monatliche Correspondenz} B. IX. S. 385\HistoricalNoteEight{Mai 1804. Band VI. S. 94.]}.
